% ===========================================================================
% CS 421 final report -- SemEval-2026 Task 4, Track A
%
% HOW THIS FILE WORKS
%   Every number in the prose is a macro from numbers.tex, which is
%   generated from reports/results/*.json and report/facts.json.
%   Never type a figure directly into a sentence. If a number has no macro,
%   add the run to the RUNS table in scripts/export_numbers.py, or add the
%   figure to report/facts.json with its source.
%
%   make numbers   regenerate numbers.tex
%   make report    regenerate numbers.tex, then build report.pdf
%   make check     flag any figure typed directly into this file
%
% FORMAT (from report_and_walkthrough.pdf)
%   PDF, MAXIMUM SIX PAGES excluding references. Font >= 11pt, letter.
%   Over-length submissions are graded only to the maximum.
%   Section order follows the six graded items:
%     1 Task and Setup 12 | 2 Data 12 | 3 Approach 22
%     4 Results 17 | 5 Analysis 27 | 6 Ethics and Limitations 10
% ===========================================================================
\documentclass[11pt,letterpaper,twocolumn]{article}
\usepackage[top=0.75in,bottom=0.75in,left=0.7in,right=0.7in]{geometry}
\setlength{\columnsep}{0.22in}
\usepackage{balance}   % even columns on the last page
\usepackage{times}
\usepackage{booktabs}
\usepackage{graphicx}
\usepackage{enumitem}
\usepackage{amsmath}
\usepackage[hidelinks]{hyperref}
\usepackage[round]{natbib}
\usepackage{titlesec}
\titlespacing*{\section}{0pt}{6pt}{1.5pt}
\titlespacing*{\subsection}{0pt}{4pt}{0.5pt}
\titleformat{\section}{\normalfont\large\bfseries}{\thesection.}{0.5em}{}
\titleformat{\subsection}{\normalfont\normalsize\bfseries}{\thesubsection.}{0.5em}{}
\setlist[enumerate]{leftmargin=1.4em,topsep=2pt,itemsep=2pt,parsep=0pt}
\setlist[itemize]{leftmargin=1.2em,topsep=2pt,itemsep=2pt,parsep=0pt}
\graphicspath{{../reports/figures/}}
% Float placement: without these, full-width floats push text onto new pages
% and the two-column gain disappears.
\renewcommand{\topfraction}{0.9}
\renewcommand{\bottomfraction}{0.7}
\renewcommand{\textfraction}{0.07}
\renewcommand{\floatpagefraction}{0.8}

% Generated by scripts/export_numbers.py -- do not edit by hand.
% Regenerate with: make numbers

% ---- split sizes --------------------------------------------------------
% DevN: labelled triples in the dev split, counted from dev__l0_jaccard__default.json
\newcommand{\DevN}{200}
% TestN: labelled triples in the test split, counted from test__l0_jaccard.json
\newcommand{\TestN}{400}
% SampleN: labelled triples in the sample split, counted from sample__l3n_characters_nocoref__e1c6459a.json
\newcommand{\SampleN}{39}

% ---- per-run figures ----------------------------------------------------
% Random: dev__l0_random__843f5ee4.json
\newcommand{\RandomAcc}{46.00}
\newcommand{\RandomN}{200}
\newcommand{\RandomLo}{39.5}
\newcommand{\RandomHi}{52.5}
\newcommand{\RandomCI}{[39.5, 52.5]}
\newcommand{\RandomSD}{49.96}
\newcommand{\RandomSE}{3.53}
\newcommand{\RandomCV}{single fit}
% Jaccard: dev__l0_jaccard__default.json
\newcommand{\JaccardAcc}{45.00}
\newcommand{\JaccardN}{200}
\newcommand{\JaccardLo}{38.0}
\newcommand{\JaccardHi}{52.0}
\newcommand{\JaccardCI}{[38.0, 52.0]}
\newcommand{\JaccardSD}{49.87}
\newcommand{\JaccardSE}{3.53}
\newcommand{\JaccardCV}{single fit}
% JaccardTest: test__l0_jaccard.json
\newcommand{\JaccardTestAcc}{58.00}
\newcommand{\JaccardTestN}{400}
\newcommand{\JaccardTestLo}{53.0}
\newcommand{\JaccardTestHi}{62.8}
\newcommand{\JaccardTestCI}{[53.0, 62.8]}
\newcommand{\JaccardTestSD}{49.42}
\newcommand{\JaccardTestSE}{2.47}
\newcommand{\JaccardTestCV}{single fit}
% Tfidf: dev__l1_tfidf__bdf0b143.json (cross-validated)
\newcommand{\TfidfAcc}{50.50}
\newcommand{\TfidfN}{200}
\newcommand{\TfidfLo}{43.5}
\newcommand{\TfidfHi}{57.5}
\newcommand{\TfidfCI}{[43.5, 57.5]}
\newcommand{\TfidfSD}{50.12}
\newcommand{\TfidfSE}{3.54}
\newcommand{\TfidfCV}{cross-validated}
% LTwo: dev__l2_e5__5d11bdc7.json
\newcommand{\LTwoAcc}{58.50}
\newcommand{\LTwoN}{200}
\newcommand{\LTwoLo}{52.0}
\newcommand{\LTwoHi}{65.5}
\newcommand{\LTwoCI}{[52.0, 65.5]}
\newcommand{\LTwoSD}{49.40}
\newcommand{\LTwoSE}{3.49}
\newcommand{\LTwoCV}{single fit}
% LThreeN: dev__l3n_characters_nocoref__e1c6459a.json
\newcommand{\LThreeNAcc}{57.00}
\newcommand{\LThreeNN}{200}
\newcommand{\LThreeNLo}{50.0}
\newcommand{\LThreeNHi}{63.5}
\newcommand{\LThreeNCI}{[50.0, 63.5]}
\newcommand{\LThreeNSD}{49.63}
\newcommand{\LThreeNSE}{3.51}
\newcommand{\LThreeNCV}{single fit}
% LThreeNActors: dev__l3n_characters_nocoref__b34df586.json
\newcommand{\LThreeNActorsAcc}{52.50}
\newcommand{\LThreeNActorsN}{200}
\newcommand{\LThreeNActorsLo}{45.5}
\newcommand{\LThreeNActorsHi}{59.5}
\newcommand{\LThreeNActorsCI}{[45.5, 59.5]}
\newcommand{\LThreeNActorsSD}{50.06}
\newcommand{\LThreeNActorsSE}{3.54}
\newcommand{\LThreeNActorsCV}{single fit}
% LThreeX: dev__l3x_characters_embed__45fe30af.json
\newcommand{\LThreeXAcc}{58.00}
\newcommand{\LThreeXN}{200}
\newcommand{\LThreeXLo}{51.0}
\newcommand{\LThreeXHi}{65.0}
\newcommand{\LThreeXCI}{[51.0, 65.0]}
\newcommand{\LThreeXSD}{49.48}
\newcommand{\LThreeXSE}{3.50}
\newcommand{\LThreeXCV}{single fit}
% LThreeXActors: dev__l3x_characters_embed__ad4ae3b7.json
\newcommand{\LThreeXActorsAcc}{58.50}
\newcommand{\LThreeXActorsN}{200}
\newcommand{\LThreeXActorsLo}{52.0}
\newcommand{\LThreeXActorsHi}{65.5}
\newcommand{\LThreeXActorsCI}{[52.0, 65.5]}
\newcommand{\LThreeXActorsSD}{49.40}
\newcommand{\LThreeXActorsSE}{3.49}
\newcommand{\LThreeXActorsCV}{single fit}
% LThreeXEFive: dev__l3x_characters_embed__ed1fb484.json
\newcommand{\LThreeXEFiveAcc}{56.00}
\newcommand{\LThreeXEFiveN}{200}
\newcommand{\LThreeXEFiveLo}{48.5}
\newcommand{\LThreeXEFiveHi}{63.0}
\newcommand{\LThreeXEFiveCI}{[48.5, 63.0]}
\newcommand{\LThreeXEFiveSD}{49.76}
\newcommand{\LThreeXEFiveSE}{3.52}
\newcommand{\LThreeXEFiveCV}{single fit}
% LFourN: dev__l4n_roles_nocoref__48c73565.json
\newcommand{\LFourNAcc}{56.00}
\newcommand{\LFourNN}{200}
\newcommand{\LFourNLo}{49.0}
\newcommand{\LFourNHi}{63.0}
\newcommand{\LFourNCI}{[49.0, 63.0]}
\newcommand{\LFourNSD}{49.76}
\newcommand{\LFourNSE}{3.52}
\newcommand{\LFourNCV}{single fit}
% LFourNActors: dev__l4n_roles_nocoref__4310ef63.json
\newcommand{\LFourNActorsAcc}{53.50}
\newcommand{\LFourNActorsN}{200}
\newcommand{\LFourNActorsLo}{47.0}
\newcommand{\LFourNActorsHi}{60.0}
\newcommand{\LFourNActorsCI}{[47.0, 60.0]}
\newcommand{\LFourNActorsSD}{50.00}
\newcommand{\LFourNActorsSE}{3.54}
\newcommand{\LFourNActorsCV}{single fit}
% LFourX: dev__l4x_roles_embed__4fcadc5a.json
\newcommand{\LFourXAcc}{57.00}
\newcommand{\LFourXN}{200}
\newcommand{\LFourXLo}{50.0}
\newcommand{\LFourXHi}{64.0}
\newcommand{\LFourXCI}{[50.0, 64.0]}
\newcommand{\LFourXSD}{49.63}
\newcommand{\LFourXSE}{3.51}
\newcommand{\LFourXCV}{single fit}
% LThree: dev__l3_characters__a3a1b934.json
\newcommand{\LThreeAcc}{50.50}
\newcommand{\LThreeN}{200}
\newcommand{\LThreeLo}{44.0}
\newcommand{\LThreeHi}{57.5}
\newcommand{\LThreeCI}{[44.0, 57.5]}
\newcommand{\LThreeSD}{50.12}
\newcommand{\LThreeSE}{3.54}
\newcommand{\LThreeCV}{single fit}
% RandomTest: test__l0_random__843f5ee4.json
\newcommand{\RandomTestAcc}{47.50}
\newcommand{\RandomTestN}{400}
\newcommand{\RandomTestLo}{42.8}
\newcommand{\RandomTestHi}{52.5}
\newcommand{\RandomTestCI}{[42.8, 52.5]}
\newcommand{\RandomTestSD}{50.00}
\newcommand{\RandomTestSE}{2.50}
\newcommand{\RandomTestCV}{single fit}
% TfidfTest: test__l1_tfidf__bdf0b143.json
\newcommand{\TfidfTestAcc}{57.50}
\newcommand{\TfidfTestN}{400}
\newcommand{\TfidfTestLo}{52.8}
\newcommand{\TfidfTestHi}{62.5}
\newcommand{\TfidfTestCI}{[52.8, 62.5]}
\newcommand{\TfidfTestSD}{49.50}
\newcommand{\TfidfTestSE}{2.47}
\newcommand{\TfidfTestCV}{single fit}
% LTwoTest: test__l2_e5__5d11bdc7.json
\newcommand{\LTwoTestAcc}{63.50}
\newcommand{\LTwoTestN}{400}
\newcommand{\LTwoTestLo}{58.8}
\newcommand{\LTwoTestHi}{68.2}
\newcommand{\LTwoTestCI}{[58.8, 68.2]}
\newcommand{\LTwoTestSD}{48.20}
\newcommand{\LTwoTestSE}{2.41}
\newcommand{\LTwoTestCV}{single fit}
% LThreeXTest: test__l3x_characters_embed__ad4ae3b7.json
\newcommand{\LThreeXTestAcc}{57.50}
\newcommand{\LThreeXTestN}{400}
\newcommand{\LThreeXTestLo}{52.5}
\newcommand{\LThreeXTestHi}{62.5}
\newcommand{\LThreeXTestCI}{[52.5, 62.5]}
\newcommand{\LThreeXTestSD}{49.50}
\newcommand{\LThreeXTestSE}{2.47}
\newcommand{\LThreeXTestCV}{single fit}
% LFiveTest: test__l5_combined__4619f43f.json
\newcommand{\LFiveTestAcc}{57.75}
\newcommand{\LFiveTestN}{400}
\newcommand{\LFiveTestLo}{53.0}
\newcommand{\LFiveTestHi}{62.5}
\newcommand{\LFiveTestCI}{[53.0, 62.5]}
\newcommand{\LFiveTestSD}{49.46}
\newcommand{\LFiveTestSE}{2.47}
\newcommand{\LFiveTestCV}{single fit}
% LFiveCorefTest: test__l5_combined__13cac587.json
\newcommand{\LFiveCorefTestAcc}{63.50}
\newcommand{\LFiveCorefTestN}{400}
\newcommand{\LFiveCorefTestLo}{58.8}
\newcommand{\LFiveCorefTestHi}{68.2}
\newcommand{\LFiveCorefTestCI}{[58.8, 68.2]}
\newcommand{\LFiveCorefTestSD}{48.20}
\newcommand{\LFiveCorefTestSE}{2.41}
\newcommand{\LFiveCorefTestCV}{single fit}
% LFour: dev__l4_roles__779c6e0c.json
\newcommand{\LFourAcc}{52.00}
\newcommand{\LFourN}{200}
\newcommand{\LFourLo}{45.0}
\newcommand{\LFourHi}{59.0}
\newcommand{\LFourCI}{[45.0, 59.0]}
\newcommand{\LFourSD}{50.09}
\newcommand{\LFourSE}{3.54}
\newcommand{\LFourCV}{single fit}
% LFourXActors: dev__l4x_roles_embed__ae3c7a38.json
\newcommand{\LFourXActorsAcc}{57.50}
\newcommand{\LFourXActorsN}{200}
\newcommand{\LFourXActorsLo}{50.5}
\newcommand{\LFourXActorsHi}{64.5}
\newcommand{\LFourXActorsCI}{[50.5, 64.5]}
\newcommand{\LFourXActorsSD}{49.56}
\newcommand{\LFourXActorsSE}{3.50}
\newcommand{\LFourXActorsCV}{single fit}
% LFiveX: dev__l5_combined__4619f43f.json (cross-validated)
\newcommand{\LFiveXAcc}{56.00}
\newcommand{\LFiveXN}{200}
\newcommand{\LFiveXLo}{49.5}
\newcommand{\LFiveXHi}{63.0}
\newcommand{\LFiveXCI}{[49.5, 63.0]}
\newcommand{\LFiveXSD}{49.76}
\newcommand{\LFiveXSE}{3.52}
\newcommand{\LFiveXCV}{cross-validated}
% LSeven: dev__l7_components__e787bd4b.json (cross-validated)
\newcommand{\LSevenAcc}{53.00}
\newcommand{\LSevenN}{200}
\newcommand{\LSevenLo}{46.5}
\newcommand{\LSevenHi}{60.0}
\newcommand{\LSevenCI}{[46.5, 60.0]}
\newcommand{\LSevenSD}{50.04}
\newcommand{\LSevenSE}{3.54}
\newcommand{\LSevenCV}{cross-validated}

% ---- paired comparisons (McNemar, same items) ---------------------------
% Jaccard 45.00 vs Tfidf 50.50  [chi-square, continuity corrected]
\newcommand{\McnJaccardVsTfidfP}{0.22}
\newcommand{\McnJaccardVsTfidfB}{28}
\newcommand{\McnJaccardVsTfidfC}{39}
\newcommand{\McnJaccardVsTfidfDisc}{67}
\newcommand{\McnJaccardVsTfidfDiff}{+5.50}
\newcommand{\McnJaccardVsTfidfTest}{chi-square, continuity corrected}
% Tfidf 50.50 vs LTwo 58.50  [chi-square, continuity corrected]
\newcommand{\McnTfidfVsLTwoP}{0.10}
\newcommand{\McnTfidfVsLTwoB}{34}
\newcommand{\McnTfidfVsLTwoC}{50}
\newcommand{\McnTfidfVsLTwoDisc}{84}
\newcommand{\McnTfidfVsLTwoDiff}{+8.00}
\newcommand{\McnTfidfVsLTwoTest}{chi-square, continuity corrected}
% LTwo 58.50 vs LThreeX 58.00  [chi-square, continuity corrected]
\newcommand{\McnLTwoVsLThreeXP}{1.00}
\newcommand{\McnLTwoVsLThreeXB}{32}
\newcommand{\McnLTwoVsLThreeXC}{31}
\newcommand{\McnLTwoVsLThreeXDisc}{63}
\newcommand{\McnLTwoVsLThreeXDiff}{-0.50}
\newcommand{\McnLTwoVsLThreeXTest}{chi-square, continuity corrected}
% LTwo 58.50 vs LThreeXActors 58.50  [chi-square, continuity corrected]
\newcommand{\McnLTwoVsLThreeXActorsP}{0.90}
\newcommand{\McnLTwoVsLThreeXActorsB}{32}
\newcommand{\McnLTwoVsLThreeXActorsC}{32}
\newcommand{\McnLTwoVsLThreeXActorsDisc}{64}
\newcommand{\McnLTwoVsLThreeXActorsDiff}{+0.00}
\newcommand{\McnLTwoVsLThreeXActorsTest}{chi-square, continuity corrected}
% LTwo 58.50 vs LFourX 57.00  [chi-square, continuity corrected]
\newcommand{\McnLTwoVsLFourXP}{0.81}
\newcommand{\McnLTwoVsLFourXB}{35}
\newcommand{\McnLTwoVsLFourXC}{32}
\newcommand{\McnLTwoVsLFourXDisc}{67}
\newcommand{\McnLTwoVsLFourXDiff}{-1.50}
\newcommand{\McnLTwoVsLFourXTest}{chi-square, continuity corrected}
% LTwo 58.50 vs LFiveX 56.00  [chi-square, continuity corrected]
\newcommand{\McnLTwoVsLFiveXP}{0.54}
\newcommand{\McnLTwoVsLFiveXB}{24}
\newcommand{\McnLTwoVsLFiveXC}{19}
\newcommand{\McnLTwoVsLFiveXDisc}{43}
\newcommand{\McnLTwoVsLFiveXDiff}{-2.50}
\newcommand{\McnLTwoVsLFiveXTest}{chi-square, continuity corrected}
% LThreeX 58.00 vs LThreeXEFive 56.00  [chi-square, continuity corrected]
\newcommand{\McnLThreeXVsLThreeXEFiveP}{0.65}
\newcommand{\McnLThreeXVsLThreeXEFiveB}{24}
\newcommand{\McnLThreeXVsLThreeXEFiveC}{20}
\newcommand{\McnLThreeXVsLThreeXEFiveDisc}{44}
\newcommand{\McnLThreeXVsLThreeXEFiveDiff}{-2.00}
\newcommand{\McnLThreeXVsLThreeXEFiveTest}{chi-square, continuity corrected}
% LThreeN 57.00 vs LThreeNActors 52.50  [chi-square, continuity corrected]
\newcommand{\McnLThreeNVsLThreeNActorsP}{0.26}
\newcommand{\McnLThreeNVsLThreeNActorsB}{30}
\newcommand{\McnLThreeNVsLThreeNActorsC}{21}
\newcommand{\McnLThreeNVsLThreeNActorsDisc}{51}
\newcommand{\McnLThreeNVsLThreeNActorsDiff}{-4.50}
\newcommand{\McnLThreeNVsLThreeNActorsTest}{chi-square, continuity corrected}
% LFourN 56.00 vs LFourNActors 53.50  [chi-square, continuity corrected]
\newcommand{\McnLFourNVsLFourNActorsP}{0.56}
\newcommand{\McnLFourNVsLFourNActorsB}{26}
\newcommand{\McnLFourNVsLFourNActorsC}{21}
\newcommand{\McnLFourNVsLFourNActorsDisc}{47}
\newcommand{\McnLFourNVsLFourNActorsDiff}{-2.50}
\newcommand{\McnLFourNVsLFourNActorsTest}{chi-square, continuity corrected}
% LThreeNActors 52.50 vs LThree 50.50  [chi-square, continuity corrected]
\newcommand{\McnLThreeNActorsVsLThreeP}{0.74}
\newcommand{\McnLThreeNActorsVsLThreeB}{44}
\newcommand{\McnLThreeNActorsVsLThreeC}{40}
\newcommand{\McnLThreeNActorsVsLThreeDisc}{84}
\newcommand{\McnLThreeNActorsVsLThreeDiff}{-2.00}
\newcommand{\McnLThreeNActorsVsLThreeTest}{chi-square, continuity corrected}
% LThreeN 57.00 vs LThree 50.50  [chi-square, continuity corrected]
\newcommand{\McnLThreeNVsLThreeP}{0.18}
\newcommand{\McnLThreeNVsLThreeB}{47}
\newcommand{\McnLThreeNVsLThreeC}{34}
\newcommand{\McnLThreeNVsLThreeDisc}{81}
\newcommand{\McnLThreeNVsLThreeDiff}{-6.50}
\newcommand{\McnLThreeNVsLThreeTest}{chi-square, continuity corrected}
% LFourNActors 53.50 vs LFour 52.00  [chi-square, continuity corrected]
\newcommand{\McnLFourNActorsVsLFourP}{0.84}
\newcommand{\McnLFourNActorsVsLFourB}{48}
\newcommand{\McnLFourNActorsVsLFourC}{45}
\newcommand{\McnLFourNActorsVsLFourDisc}{93}
\newcommand{\McnLFourNActorsVsLFourDiff}{-1.50}
\newcommand{\McnLFourNActorsVsLFourTest}{chi-square, continuity corrected}
% LFourN 56.00 vs LFour 52.00  [chi-square, continuity corrected]
\newcommand{\McnLFourNVsLFourP}{0.46}
\newcommand{\McnLFourNVsLFourB}{49}
\newcommand{\McnLFourNVsLFourC}{41}
\newcommand{\McnLFourNVsLFourDisc}{90}
\newcommand{\McnLFourNVsLFourDiff}{-4.00}
\newcommand{\McnLFourNVsLFourTest}{chi-square, continuity corrected}
% LTwoTest 63.50 vs LThreeXTest 57.50  [chi-square, continuity corrected]
\newcommand{\McnLTwoTestVsLThreeXTestP}{0.033}
\newcommand{\McnLTwoTestVsLThreeXTestB}{70}
\newcommand{\McnLTwoTestVsLThreeXTestC}{46}
\newcommand{\McnLTwoTestVsLThreeXTestDisc}{116}
\newcommand{\McnLTwoTestVsLThreeXTestDiff}{-6.00}
\newcommand{\McnLTwoTestVsLThreeXTestTest}{chi-square, continuity corrected}
% LTwoTest 63.50 vs LFiveTest 57.75  [chi-square, continuity corrected]
\newcommand{\McnLTwoTestVsLFiveTestP}{0.027}
\newcommand{\McnLTwoTestVsLFiveTestB}{61}
\newcommand{\McnLTwoTestVsLFiveTestC}{38}
\newcommand{\McnLTwoTestVsLFiveTestDisc}{99}
\newcommand{\McnLTwoTestVsLFiveTestDiff}{-5.75}
\newcommand{\McnLTwoTestVsLFiveTestTest}{chi-square, continuity corrected}
% LTwoTest 63.50 vs TfidfTest 57.50  [chi-square, continuity corrected]
\newcommand{\McnLTwoTestVsTfidfTestP}{0.085}
\newcommand{\McnLTwoTestVsTfidfTestB}{101}
\newcommand{\McnLTwoTestVsTfidfTestC}{77}
\newcommand{\McnLTwoTestVsTfidfTestDisc}{178}
\newcommand{\McnLTwoTestVsTfidfTestDiff}{-6.00}
\newcommand{\McnLTwoTestVsTfidfTestTest}{chi-square, continuity corrected}

% ---- error correlation against Tfidf ------------------------------------
% PhiLTwo: error-vector correlation of dev__l2_e5__5d11bdc7.json with dev__l1_tfidf__bdf0b143.json, same 200 dev items
\newcommand{\PhiLTwo}{+0.16}
% PhiLThreeN: error-vector correlation of dev__l3n_characters_nocoref__e1c6459a.json with dev__l1_tfidf__bdf0b143.json, same 200 dev items
\newcommand{\PhiLThreeN}{+0.29}
% PhiLThreeX: error-vector correlation of dev__l3x_characters_embed__45fe30af.json with dev__l1_tfidf__bdf0b143.json, same 200 dev items
\newcommand{\PhiLThreeX}{+0.25}
% PhiLThreeXEFive: error-vector correlation of dev__l3x_characters_embed__ed1fb484.json with dev__l1_tfidf__bdf0b143.json, same 200 dev items
\newcommand{\PhiLThreeXEFive}{+0.27}
% PhiLFourN: error-vector correlation of dev__l4n_roles_nocoref__48c73565.json with dev__l1_tfidf__bdf0b143.json, same 200 dev items
\newcommand{\PhiLFourN}{+0.33}
% PhiLFourX: error-vector correlation of dev__l4x_roles_embed__4fcadc5a.json with dev__l1_tfidf__bdf0b143.json, same 200 dev items
\newcommand{\PhiLFourX}{+0.27}
% PhiLFiveX: error-vector correlation of dev__l5_combined__4619f43f.json with dev__l1_tfidf__bdf0b143.json, same 200 dev items
\newcommand{\PhiLFiveX}{+0.21}
% PhiMin/PhiMax: range of the 7 correlations listed above
\newcommand{\PhiMin}{+0.16}
\newcommand{\PhiMax}{+0.33}

% ---- combination ceiling ------------------------------------------------
% pool: l2_e5, l3x_characters_embed, l4x_roles_embed, l3n_characters_nocoref, l4n_roles_nocoref
\newcommand{\PoolNames}{LTwo, LThreeX, LFourX, LThreeN, LFourN}
\newcommand{\PoolSize}{5}
\newcommand{\OracleUnion}{86.0}
\newcommand{\OracleIndep}{98.6}
\newcommand{\MajorityVote}{60.0}

% ---- fitted weights -----------------------------------------------------
% LFiveX mean weights over 5 folds: l2_e5=0.32, l3x_characters_embed=0.18, l4x_roles_embed=0.50
\newcommand{\LFiveXFolds}{5}
\newcommand{\LFiveXWeights}{0.32, 0.18, 0.50}

% ---- unanimous strata, outcomes -----------------------------------------
% TfidfMarked: dev__l1_tfidf__bdf0b143.json, unanimous stratum where every annotator marked outcomes
\newcommand{\TfidfMarkedAcc}{57.5}
\newcommand{\TfidfMarkedN}{40}
% TfidfUnmarked: dev__l1_tfidf__bdf0b143.json, unanimous stratum where no annotator marked outcomes
\newcommand{\TfidfUnmarkedAcc}{53.5}
\newcommand{\TfidfUnmarkedN}{71}
% LTwoMarked: dev__l2_e5__5d11bdc7.json, unanimous stratum where every annotator marked outcomes
\newcommand{\LTwoMarkedAcc}{50.0}
\newcommand{\LTwoMarkedN}{40}
% LTwoUnmarked: dev__l2_e5__5d11bdc7.json, unanimous stratum where no annotator marked outcomes
\newcommand{\LTwoUnmarkedAcc}{67.6}
\newcommand{\LTwoUnmarkedN}{71}
% LThreeXMarked: dev__l3x_characters_embed__45fe30af.json, unanimous stratum where every annotator marked outcomes
\newcommand{\LThreeXMarkedAcc}{67.5}
\newcommand{\LThreeXMarkedN}{40}
% LThreeXUnmarked: dev__l3x_characters_embed__45fe30af.json, unanimous stratum where no annotator marked outcomes
\newcommand{\LThreeXUnmarkedAcc}{57.7}
\newcommand{\LThreeXUnmarkedN}{71}

% ---- component validity (generated) -------------------------------------
% from components__dev.json, split dev
\newcommand{\CompThemeAcc}{55.0}
\newcommand{\CompThemeCI}{[48.1, 61.7]}
\newcommand{\CompThemeLo}{48.1}
\newcommand{\CompThemeHi}{61.7}
%   theme: interval contains 50
\newcommand{\CompActionAcc}{54.5}
\newcommand{\CompActionCI}{[47.6, 61.3]}
\newcommand{\CompActionLo}{47.6}
\newcommand{\CompActionHi}{61.3}
%   action: interval contains 50
\newcommand{\CompOutcomeAcc}{54.0}
\newcommand{\CompOutcomeCI}{[47.1, 60.8]}
\newcommand{\CompOutcomeLo}{47.1}
\newcommand{\CompOutcomeHi}{60.8}
%   outcome: interval contains 50
\newcommand{\CompN}{200}

% ---- triple geometry per space (generated) -------------------------------
% from geometry__dev.json, split dev
% Tfidf: TF-IDF, unit-normalized
\newcommand{\GeomTfidfRight}{1.394}
\newcommand{\GeomTfidfWrong}{1.395}
\newcommand{\GeomTfidfDiff}{-0.0013}
\newcommand{\GeomTfidfCosRight}{+0.028}
\newcommand{\GeomTfidfCosWrong}{+0.027}
\newcommand{\GeomTfidfCloserK}{104}
\newcommand{\GeomTfidfCloserPct}{52.0}
\newcommand{\GeomTfidfCloserP}{0.62}
\newcommand{\GeomTfidfAngle}{60.5}
\newcommand{\GeomTfidfAngleNull}{60.9}
\newcommand{\GeomTfidfN}{200}
% MiniLM: all-MiniLM-L6-v2
\newcommand{\GeomMiniLMRight}{1.173}
\newcommand{\GeomMiniLMWrong}{1.188}
\newcommand{\GeomMiniLMDiff}{-0.0148}
\newcommand{\GeomMiniLMCosRight}{+0.312}
\newcommand{\GeomMiniLMCosWrong}{+0.294}
\newcommand{\GeomMiniLMCloserK}{110}
\newcommand{\GeomMiniLMCloserPct}{55.0}
\newcommand{\GeomMiniLMCloserP}{0.18}
\newcommand{\GeomMiniLMAngle}{61.3}
\newcommand{\GeomMiniLMAngleNull}{73.5}
\newcommand{\GeomMiniLMN}{200}
% EFive: e5-base-v2
\newcommand{\GeomEFiveRight}{0.638}
\newcommand{\GeomEFiveWrong}{0.647}
\newcommand{\GeomEFiveDiff}{-0.0086}
\newcommand{\GeomEFiveCosRight}{+0.796}
\newcommand{\GeomEFiveCosWrong}{+0.791}
\newcommand{\GeomEFiveCloserK}{117}
\newcommand{\GeomEFiveCloserPct}{58.5}
\newcommand{\GeomEFiveCloserP}{0.019}
\newcommand{\GeomEFiveAngle}{61.4}
\newcommand{\GeomEFiveAngleNull}{undefined}
\newcommand{\GeomEFiveN}{200}

% ---- lexical direction (generated) --------------------------------------
% all_tokens, dev, n=200
\newcommand{\LexAllJSPct}{49.5}
\newcommand{\LexAllJSK}{99}
\newcommand{\LexAllJSP}{0.94}
\newcommand{\LexAllKLAnchorPct}{42.5}
\newcommand{\LexAllKLAnchorK}{85}
\newcommand{\LexAllKLAnchorP}{0.04}
\newcommand{\LexAllKLCandPct}{56.0}
\newcommand{\LexAllKLCandK}{112}
\newcommand{\LexAllKLCandP}{0.10}
\newcommand{\LexAllCrossEntPct}{45.5}
\newcommand{\LexAllCrossEntK}{91}
\newcommand{\LexAllCrossEntP}{0.23}
\newcommand{\LexAllEditPct}{44.0}
\newcommand{\LexAllEditK}{88}
\newcommand{\LexAllEditP}{0.10}
\newcommand{\LexAllEditNormPct}{51.5}
\newcommand{\LexAllEditNormK}{103}
\newcommand{\LexAllEditNormP}{0.72}
\newcommand{\LexAllLengthPct}{47.0}
\newcommand{\LexAllLengthK}{94}
\newcommand{\LexAllLengthP}{0.44}
\newcommand{\LexAllTypeTokenPct}{52.0}
\newcommand{\LexAllTypeTokenK}{104}
\newcommand{\LexAllTypeTokenP}{0.62}
% content_words, dev, n=200
\newcommand{\LexContentJSPct}{49.0}
\newcommand{\LexContentJSK}{98}
\newcommand{\LexContentJSP}{0.83}
\newcommand{\LexContentKLAnchorPct}{43.0}
\newcommand{\LexContentKLAnchorK}{86}
\newcommand{\LexContentKLAnchorP}{0.056}
\newcommand{\LexContentKLCandPct}{55.0}
\newcommand{\LexContentKLCandK}{110}
\newcommand{\LexContentKLCandP}{0.18}
\newcommand{\LexContentCrossEntPct}{40.5}
\newcommand{\LexContentCrossEntK}{81}
\newcommand{\LexContentCrossEntP}{0.009}
\newcommand{\LexContentEditPct}{36.0}
\newcommand{\LexContentEditK}{72}
\newcommand{\LexContentEditP}{<0.001}
\newcommand{\LexContentEditNormPct}{37.5}
\newcommand{\LexContentEditNormK}{75}
\newcommand{\LexContentEditNormP}{<0.001}
\newcommand{\LexContentLengthPct}{47.5}
\newcommand{\LexContentLengthK}{95}
\newcommand{\LexContentLengthP}{0.52}
\newcommand{\LexContentTypeTokenPct}{51.5}
\newcommand{\LexContentTypeTokenK}{103}
\newcommand{\LexContentTypeTokenP}{0.72}
\newcommand{\LexN}{200}

% ---- spread across configurations (generated) ---------------------------
% computed here over the 20 distinct dev configurations in reports/results
\newcommand{\RunCountDev}{20}
\newcommand{\SpreadMeanDev}{54.05}
\newcommand{\SpreadDevSD}{3.82}
\newcommand{\SpreadMinDev}{45.00}
\newcommand{\SpreadMaxDev}{58.50}
% computed here over the 7 distinct test configurations in reports/results
\newcommand{\RunCountTest}{7}
\newcommand{\SpreadMeanTest}{57.89}
\newcommand{\SpreadTestSD}{5.34}
\newcommand{\SpreadMinTest}{47.50}
\newcommand{\SpreadMaxTest}{63.50}

% ---- tokenizer reconciliation (generated) -------------------------------
% scripts/tokenizer_reconciliation.py: the L0 token-overlap baseline on test under 10 tokenizations, against the published 56.25\%
\newcommand{\JaccardGapPts}{1.75}
\newcommand{\JaccardGapItems}{7}
\newcommand{\JaccardFlipN}{17}
\newcommand{\JaccardFlipFor}{12}
\newcommand{\JaccardFlipAgainst}{5}
\newcommand{\JaccardDiffTypes}{930}
\newcommand{\JaccardDiffApostPct}{74.7}
\newcommand{\JaccardVariantsMatching}{4}
% scripts/tokenizer_reconciliation.py: the same 10 tokenizations on dev, showing the below-chance result is not an artifact of the choice
\newcommand{\JaccardMinVariant}{44.00}
\newcommand{\JaccardMaxVariant}{47.50}
\newcommand{\JaccardVariantsN}{10}

% ---- notebook-derived figures (report/facts.json) -----------------------
% GeomRightMedian: notebooks/03_eda_geometry.ipynb, cosine-distance cell
\newcommand{\GeomRightMedian}{1.390}
% GeomWrongMedian: notebooks/03_eda_geometry.ipynb, cosine-distance cell
\newcommand{\GeomWrongMedian}{1.391}
% GeomRightCloser: notebooks/03_eda_geometry.ipynb cell 10, 76/150
\newcommand{\GeomRightCloser}{51}
% GeomAngle: notebooks/03_eda_geometry.ipynb cell 10
\newcommand{\GeomAngle}{60.6}
% AnnThemePct: majority of annotators marked the dimension; dev and test pooled, n=600. Verified by recomputation 17 Nov: 590/317/144 of 600 exactly
\newcommand{\AnnThemePct}{98.3}
% AnnActionPct: majority of annotators marked the dimension; dev and test pooled, n=600. Verified by recomputation 17 Nov: 590/317/144 of 600 exactly
\newcommand{\AnnActionPct}{52.8}
% AnnOutcomePct: majority of annotators marked the dimension; dev and test pooled, n=600. Verified by recomputation 17 Nov: 590/317/144 of 600 exactly
\newcommand{\AnnOutcomePct}{24.0}
% AnnCommentCount: notebooks/03_eda_geometry.ipynb, free-text comment count
\newcommand{\AnnCommentCount}{226}
% AnnTripleCount: sample + dev + test triples carrying annotator metadata
\newcommand{\AnnTripleCount}{600}
% AnnProseThemePct: notebooks/03_eda_geometry.ipynb, keyword scan of free-text comments
\newcommand{\AnnProseThemePct}{80.1}
% AnnProseCharPct: notebooks/03_eda_geometry.ipynb, keyword scan of free-text comments
\newcommand{\AnnProseCharPct}{33.6}
% AnnProseSettingPct: notebooks/03_eda_geometry.ipynb, keyword scan of free-text comments
\newcommand{\AnnProseSettingPct}{10.6}
% AnnCommentWords: notebooks/03_eda_geometry.ipynb, mean comment length
\newcommand{\AnnCommentWords}{38}
% AnnStoryWords: notebooks/03_eda_geometry.ipynb, mean summary length
\newcommand{\AnnStoryWords}{116}
% JaccardMargin: reports/results dev l0_jaccard margins, mean signed margin
\newcommand{\JaccardMargin}{-0.0024}
% PowerStrataMDE: scripts/power_analysis.py, between-strata contrast
\newcommand{\PowerStrataMDE}{17.3}
% PowerStrataSmall: scripts/power_analysis.py, unanimous outcomes strata on test
\newcommand{\PowerStrataSmall}{79}
% PowerStrataLarge: scripts/power_analysis.py, unanimous outcomes strata on test
\newcommand{\PowerStrataLarge}{133}
% PowerPairedLo: scripts/power_analysis.py, paired McNemar on 400 test items
\newcommand{\PowerPairedLo}{4.3}
% PowerPairedHi: scripts/power_analysis.py, paired McNemar on 400 test items
\newcommand{\PowerPairedHi}{6.2}
% OfficialBestA: Hatzel et al. 2026 task overview, Table 2 (COGNAC)
\newcommand{\OfficialBestA}{78.00}
% OfficialBestB: Hatzel et al. 2026 task overview, Table 3 (COGNAC)
\newcommand{\OfficialBestB}{72.00}
% OfficialBaselineA: Hatzel et al. 2026 task overview, zero-shot GPT-4o-mini baseline
\newcommand{\OfficialBaselineA}{67.00}
% OfficialMeanA: computed here from the ranked Track A systems in Hatzel et al. 2026
\newcommand{\OfficialMeanA}{66.51}
% OfficialRankedA: Hatzel et al. 2026 task overview, ranked Track A systems
\newcommand{\OfficialRankedA}{44}
% OfficialTeams: Hatzel et al. 2026 task overview
\newcommand{\OfficialTeams}{46}
% OfficialRuns: Hatzel et al. 2026 task overview
\newcommand{\OfficialRuns}{71}
% HumanBound: Hatzel et al. 2026 task overview, estimated human upper bound
\newcommand{\HumanBound}{89}
% AlphaOverall: Hatzel et al. 2026 task overview, Krippendorff alpha
\newcommand{\AlphaOverall}{0.33}
% AlphaPerDimLo: Hatzel et al. 2026 task overview, per-dimension Krippendorff alpha
\newcommand{\AlphaPerDimLo}{0.05}
% AlphaPerDimHi: Hatzel et al. 2026 task overview, per-dimension Krippendorff alpha
\newcommand{\AlphaPerDimHi}{0.11}
% StoryNetAcc: Vijayachandran et al. 2026 (StoryNet) system paper
\newcommand{\StoryNetAcc}{57.75}
% StoryNetRank: Hatzel et al. 2026 task overview, Track A ranking
\newcommand{\StoryNetRank}{39}
% PublishedJaccardTest: Hatzel et al. 2026 task overview, lexical baseline on test
\newcommand{\PublishedJaccardTest}{56.25}
% PublishedTfidfTest: Hatzel et al. 2026 task overview, lexical baseline on test
\newcommand{\PublishedTfidfTest}{58.25}
% ChaturvediBefore: Chaturvedi et al. 2018, precision at rank 1, event model
\newcommand{\ChaturvediBefore}{0.558}
% ChaturvediAfter: Chaturvedi et al. 2018, precision at rank 1, character model
\newcommand{\ChaturvediAfter}{0.637}
% BarancovaGain: Barancova and Sileno 2026, relative retrieval improvement, per cent
\newcommand{\BarancovaGain}{14.6}
% ElsnerSentimentOnly: Elsner 2012, sentiment-trajectory-only accuracy
\newcommand{\ElsnerSentimentOnly}{50}
% GeomN: notebooks/03_eda_geometry.ipynb cell 10
\newcommand{\GeomN}{150}
% GeomRightCloserK: notebooks/03_eda_geometry.ipynb cell 10
\newcommand{\GeomRightCloserK}{76}
% GeomRightCloserP: exact two-sided binomial on 76/150 against 0.5, computed in scripts
\newcommand{\GeomRightCloserP}{0.93}
% GeomDiff: notebooks/03_eda_geometry.ipynb cell 10
\newcommand{\GeomDiff}{-0.0008}
% LFiveNInSample: reports/archive/dev__l5_combined__over-L3n__IN-SAMPLE-NOT-CV.json
\newcommand{\LFiveNInSample}{59.00}
% StratThemeMarked: partition of the dev split on 'abstract_theme', annotator labels only (no model)
\newcommand{\StratThemeMarked}{190}
% StratThemeUnmarked: partition of the dev split on 'abstract_theme', annotator labels only (no model)
\newcommand{\StratThemeUnmarked}{0}
% StratActionMarked: partition of the dev split on 'course_of_actions', annotator labels only (no model)
\newcommand{\StratActionMarked}{88}
% StratActionUnmarked: partition of the dev split on 'course_of_actions', annotator labels only (no model)
\newcommand{\StratActionUnmarked}{19}
% StratOutcomeMarked: partition of the dev split on 'outcomes', annotator labels only (no model)
\newcommand{\StratOutcomeMarked}{40}
% StratOutcomeUnmarked: partition of the dev split on 'outcomes', annotator labels only (no model)
\newcommand{\StratOutcomeUnmarked}{71}
% AnnDimsPerJudgment: computed from the dev split: mean number of the three dimensions ticked per annotator judgment
\newcommand{\AnnDimsPerJudgment}{2.08}
% AnnThemeJudgPct: computed from the dev split: share of individual annotator judgments ticking this dimension (434 judgments)
\newcommand{\AnnThemeJudgPct}{97.7}
% AnnActionJudgPct: computed from the dev split: share of individual annotator judgments ticking this dimension (434 judgments)
\newcommand{\AnnActionJudgPct}{68.2}
% AnnOutcomeJudgPct: computed from the dev split: share of individual annotator judgments ticking this dimension (434 judgments)
\newcommand{\AnnOutcomeJudgPct}{41.7}
% AnnPoolN: dev (200) + test (400) triples, the pool the majority-rule rates are computed over
\newcommand{\AnnPoolN}{600}
% AnnotTotal: Hatzel et al. 2026 task overview, annotation pipeline
\newcommand{\AnnotTotal}{1,237}
% AnnotAgreed: Hatzel et al. 2026 task overview, annotation pipeline
\newcommand{\AnnotAgreed}{831}
% AnnotAdjudicated: Hatzel et al. 2026 task overview, annotation pipeline
\newcommand{\AnnotAdjudicated}{208}
% AnnotKept: Hatzel et al. 2026 task overview, annotation pipeline
\newcommand{\AnnotKept}{1,039}
% AnnotAgreePct: Hatzel et al. 2026 task overview, annotation pipeline
\newcommand{\AnnotAgreePct}{67}
% AnnotDropped: derived, not reported: Hatzel et al. 2026 give 1,237 annotated and 1,039 kept, so 198 were not carried forward
\newcommand{\AnnotDropped}{198}
% PhiNullLo: permutation test computed here: one layer's per-item error vector shuffled 4,000 times with its accuracy held fixed, phi recomputed each time
\newcommand{\PhiNullLo}{-0.14}
% PhiNullHi: permutation test computed here: one layer's per-item error vector shuffled 4,000 times with its accuracy held fixed, phi recomputed each time
\newcommand{\PhiNullHi}{+0.14}
% PhiNullPLTwo: permutation test computed here: one layer's per-item error vector shuffled 4,000 times with its accuracy held fixed, phi recomputed each time; L2 against TF-IDF
\newcommand{\PhiNullPLTwo}{0.021}
% PhiNullPLThreeX: permutation test computed here: one layer's per-item error vector shuffled 4,000 times with its accuracy held fixed, phi recomputed each time; L3x against TF-IDF
\newcommand{\PhiNullPLThreeX}{0.001}
% TfidfWrongN: computed: items the TF-IDF baseline gets wrong on dev
\newcommand{\TfidfWrongN}{99}
% RecoverLTwo: computed: of the 99 items TF-IDF gets wrong on dev, the number this layer gets right
\newcommand{\RecoverLTwo}{50}
% RecoverLThreeN: computed: of the 99 items TF-IDF gets wrong on dev, the number this layer gets right
\newcommand{\RecoverLThreeN}{42}
% DistractorPercentile: Hatzel et al. 2026 task overview: the similarity percentile band the distractor is drawn from
\newcommand{\DistractorPercentile}{98th to 99th}
% ExampleCastFull: scripts/examples.py, development triple 0: noun-chunk heads before filtering
\newcommand{\ExampleCastFull}{59}
% ExampleCastActors: scripts/examples.py, development triple 0: heads surviving the actors-only filter
\newcommand{\ExampleCastActors}{21}
% WeightRangeLTwo: reports/results/dev__l5_combined: per-fold weight on l2_e5 across 5 cross-validation folds
\newcommand{\WeightRangeLTwo}{0.2 to 0.6}
% WeightRangeLThreeX: reports/results/dev__l5_combined: per-fold weight on l3x_characters_embed across 5 cross-validation folds
\newcommand{\WeightRangeLThreeX}{0.0 to 0.4}
% WeightRangeLFourX: reports/results/dev__l5_combined: per-fold weight on l4x_roles_embed across 5 cross-validation folds
\newcommand{\WeightRangeLFourX}{0.0 to 0.7}
% BestSingleAcc: computed: highest dev accuracy of any single layer (l2_e5 and l3x actors both reach it)
\newcommand{\BestSingleAcc}{58.50}
% CombinationLoss: computed: best single layer minus the cross-validated combination
\newcommand{\CombinationLoss}{2.50}
% OracleShortfall: computed: independence expectation minus observed union, five-layer pool
\newcommand{\OracleShortfall}{12.6}
% OracleResidue: computed: per cent of dev items no layer in the five-layer pool gets right
\newcommand{\OracleResidue}{14.0}
% LexTestCrossEntPct: scripts/lexical_direction.py on test, content words, cross_entropy
\newcommand{\LexTestCrossEntPct}{49.0}
% LexTestCrossEntP: scripts/lexical_direction.py on test, content words, cross_entropy
\newcommand{\LexTestCrossEntP}{0.73}
% LexTestEditPct: scripts/lexical_direction.py on test, content words, edit_distance
\newcommand{\LexTestEditPct}{45.8}
% LexTestEditP: scripts/lexical_direction.py on test, content words, edit_distance
\newcommand{\LexTestEditP}{0.10}
% FreezeSha: git rev-parse --short HEAD at the freeze
\newcommand{\FreezeSha}{a9b9881}
% LFiveCorefWeights: reports/results test l5_combined over the coreference variants: weights fitted on dev
\newcommand{\LFiveCorefWeights}{1.0, 0.0, 0.0}
% LFiveTestWeights: reports/results test l5_combined over L3x/L4x: weights fitted on dev
\newcommand{\LFiveTestWeights}{0.3, 0.0, 0.7}
% SrlArgMatchLLM: Bonn et al. 2024: SRL by LLM prompt, argument match
\newcommand{\SrlArgMatchLLM}{48.6}
% SrlArgMatchHuman: Bonn et al. 2024: human argument match, same evaluation
\newcommand{\SrlArgMatchHuman}{88.3}
% GrammarWeight: CascadeMind system description via Hatzel et al. 2026: learned weight on the five-phase story grammar
\newcommand{\GrammarWeight}{0.40}
% TfidfWeight: CascadeMind system description via Hatzel et al. 2026: learned weight on TF-IDF
\newcommand{\TfidfWeight}{0.49}
% EmbedWeight: CascadeMind system description via Hatzel et al. 2026: learned weight on embeddings
\newcommand{\EmbedWeight}{0.08}
% ArcWeight: CascadeMind system description via Hatzel et al. 2026: learned weight on the tension curve
\newcommand{\ArcWeight}{0.02}
% SegPadPct: computed from the 600 dev summaries: share with fewer than 5 sentences, which split_segments pads with a repeated sentence
\newcommand{\SegPadPct}{26}
% SegMedianSents: computed from the 600 dev summaries: median sentence count
\newcommand{\SegMedianSents}{6}
% ExampleJaccardRight: scripts/examples.py, development triple 160: Jaccard against the anchor over the L0 tokenizer (candidate chosen by annotators)
\newcommand{\ExampleJaccardRight}{0.0763}
% ExampleJaccardWrong: scripts/examples.py, development triple 160: Jaccard against the anchor over the L0 tokenizer (distractor)
\newcommand{\ExampleJaccardWrong}{0.1515}
% ExampleJaccardGap: scripts/examples.py, development triple 160: Jaccard against the anchor over the L0 tokenizer (margin by which overlap prefers the distractor)
\newcommand{\ExampleJaccardGap}{0.0752}
% ExampleCastRemoved: scripts/examples.py, development triple 0: noun heads discarded by the actors-only filter
\newcommand{\ExampleCastRemoved}{38}
% TrackAKept: Hatzel et al. 2026: the 1,039 kept triples span both tracks; Track A receives 39 sample + 200 dev + 400 test
\newcommand{\TrackAKept}{639}

% MacroCount: counted here, the number of macros in this file
\newcommand{\MacroCount}{534}



\title{\vspace{-2em}\textbf{Does Narrative Structure Help Where Words Do Not?}\\[2pt]
\normalsize A Layered Ablation on SemEval-2026 Task 4, Track A}
\author{Claude Kwizera \\ \small CS 421: Natural Language Processing \\
\small University of Illinois Chicago}
\date{\small November 25, 2026}

\begin{document}
\twocolumn[
\begin{@twocolumnfalse}
\maketitle

% ---------------------------------------------------------------------------
% ABSTRACT
%   CLAIM    Structure did not add signal, and the reason is measurable.
%   BUDGET   140-170 words. Write this last; revise it first.
% ---------------------------------------------------------------------------
\begin{abstract}
\noindent
SemEval-2026 Task~4 asks which of two story summaries is more similar to an  % ok
anchor. Because the triples that survive annotation are selected to defeat
surface overlap, a Jaccard baseline scores \JaccardAcc\% on development data,
below chance. This report asks whether narrative structure supplies the signal
that wording does not, through an ablation in which each layer adds one
representation. No layer separates from any other on development data, and on the held-out
split the embedding backbone leads the best structural layer,
\LTwoTestAcc\% against \LThreeXTestAcc\%
($p = \McnLTwoTestVsLThreeXTestP$). Three convergent diagnostics explain why.
The alignment compares spans inside an embedding space in which even whole
summaries fail to separate the candidates; the layers err on the same items as
the lexical baseline; and their combination ceiling falls \OracleShortfall\
points below what independent systems would reach.
\end{abstract}
\vspace{0.4em}
\end{@twocolumnfalse}
]

% ===========================================================================
% READ THIS BEFORE WRITING A WORD
%
% The generated numbers do not say what an earlier reading of the table
% suggested. Two corrections:
%
%  1. NOTHING IS SIGNIFICANT ON DEV. The words-to-meaning jump is +8.00
%     points but p = 0.10, from 84 discordant pairs. Do not write "the only
%     significant difference is L1 to L2" -- there is no significant
%     difference. The correct sentence is that the ablation separates nothing
%     at n = 200, which is exactly what the power analysis predicted, and
%     that this is why Section 7 does not rest on accuracy at all.
%
%  2. THE UNION CEILING DEPENDS ON THE POOL. \OracleUnion over \PoolSize
%     layers is not the same quantity as any union figure computed over three
%     layers earlier in the project. The pool is named in \PoolNames and in
%     POOL in scripts/export_numbers.py. Quote the pool whenever you quote
%     the number.
%
% Both corrections make the report stronger: the argument moves from "our
% layer lost by a little" to "accuracy cannot settle this, so here is what
% can", which is the section you actually have evidence for.
% ===========================================================================

% ---------------------------------------------------------------------------

% ===========================================================================
\section{Task and Setup}\label{sec:task}
% RUBRIC 1 (12 pts). What is the task, what does it demand, what was built.
% No construction mechanism here -- that is Section 2.
% ===========================================================================
This project addresses Track~A of SemEval-2026 Task~4 \citep{hatzel2026}, which
asks a system to select, from two candidate story summaries, the one more
similar to an anchor.

The triples are not drawn at random from a corpus. Both candidates are
retrieved into the anchor's immediate neighbourhood by a narrative encoder, and
only those triples that two commercial language models disagree on survive to
be labelled. The released items are therefore, by construction, those on which
the obvious methods already fail. Because that construction explains almost
every result reported here, the pipeline is set out in full in
Section~\ref{sec:data} before any system is described.

The task is a binary forced choice, so the official metric is accuracy and
chance is 50\%.  % ok
Each item scores 0 or 1, which means the variance and standard deviation are
fixed once the accuracy is known and add nothing to it. Uncertainty is
therefore described with bootstrap and Wilson intervals, and systems are
compared with paired tests on the same items. On that scale the best official
system (COGNAC) reached \OfficialBestA\%, the strongest organiser baseline, a
zero-shot \texttt{GPT-4o-mini} prompt, reached \OfficialBaselineA\%, and the
estimated human upper bound is \HumanBound\%. The mean of the
\OfficialRankedA\ ranked Track~A systems, computed here rather than published,
is \OfficialMeanA\%. Across the \RunCountDev\ development configurations built
for this project the standard deviation of accuracy is \SpreadDevSD\ points,
which is about the size of the smallest difference this design can detect.

\paragraph{Setup.} The system is a sequence of layers, each adding exactly one
representational commitment to the one before it, so any change in accuracy is
attributable to that single addition and each result --- positive or negative
--- determines whether the next layer is worth building. All models are
pretrained and frozen. Parameters are fitted by cross-validation, never on the
split being scored. Because the power analysis gives a minimum detectable paired
difference of \PowerPairedLo\ to \PowerPairedHi\ points on test, one
confirmatory comparison was fixed in advance in a written pre-registration; everything else is
labelled descriptive or exploratory.

% ===========================================================================
\section{Data}\label{sec:data}
% RUBRIC 2 (12 pts). What data, how much, and how it differs from the paper.
% ===========================================================================
The official dataset is publicly
released.\footnote{\raggedright\url{https://github.com/narrative-similarity-task/narrative-similarity-dataset}}
Track~A provides \SampleN\ sample, \DevN\ development and \TestN\ test triples.
Summaries are licensed CC BY-SA 4.0 and the labels CC0:  % ok the summaries may be
reused with attribution and under the same licence, while the annotations carry
no restriction.

Three properties of the release shaped this project. The first is that no training
split is provided, so fine-tuning is unavailable: every model is pretrained and
frozen, and all engineering happens in the representation, not the weights. The second is that the test labels
are public, which makes the leaderboard reproducible but also makes
overfitting it a matter of discipline rather than of access --- one team was
disqualified on exactly this point \citep{hatzel2026}. I therefore reserved the
split for a single final evaluation, fixed in advance, having used it once
earlier for untuned baselines while the proposal was written. The third is that stories
recur across splits, so a development-selected configuration is optimistic on
test for that reason alone; fitting is therefore cross-validated within a
split.

\subsection{How the triples were built}


Summaries come from \emph{Tell Me Again}, restricted to four to eight
sentences. Each anchor is sampled at random, and both candidates are then drawn
relative to it with \texttt{story-emb}, an encoder trained so that different
summaries of the same story land close together. The nearest summary becomes
the assumed answer, while the distractor is taken from the
\DistractorPercentile\ similarity percentile, so that both candidates sit inside
the anchor's neighbourhood rather than merely near it. Two commercial language
models, which the overview does not name, then judge every triple, and only
those they disagree on reach human annotation. Of \AnnotTotal\ triples,
\AnnotAgreed\ had two annotators agree and \AnnotAdjudicated\ were settled by a
third, leaving \AnnotKept. Those \AnnotKept\ span both tracks, of which Track~A
receives \TrackAKept: the \SampleN, \DevN\ and \TestN\ splits above.

\subsection{What the annotators recorded}


Alongside the binary choice, annotators recorded which of three dimensions
applied: abstract theme, course of action, and outcomes. Because these are
independent checkboxes rather than a ranking, an annotator ticks
\AnnDimsPerJudgment\ of the three on average, and across the \AnnPoolN\
annotated triples a majority marked theme on \AnnThemePct\%, course of action
on \AnnActionPct\% and outcomes on \AnnOutcomePct\%.

Agreement on \emph{why} is far weaker than on \emph{which}: Krippendorff's
$\alpha$ is \AlphaOverall\ overall and between \AlphaPerDimLo\ and
\AlphaPerDimHi\ per dimension. That caps what any per-dimension diagnostic can
establish, and it has a further consequence, since a dimension recorded on
almost every item has no contrast group at all. Theme is marked unanimously on
\StratThemeMarked\ of \DevN\ development triples and unanimously declined on
\StratThemeUnmarked, against \StratOutcomeMarked\ and \StratOutcomeUnmarked\
for outcomes. The most cited dimension is therefore the least usable, and
outcomes, the least cited, is the only one with strata large enough to
compare.

Construction and annotation together describe a benchmark on which surface
methods should not be expected to succeed. Why that is so, rather than merely
that it is, is the subject of Section~\ref{sec:analysis}.

% ---------------------------------------------------------------------------

\subsection{Reconciliation against the published baseline}
Before trusting any number from this pipeline I reproduced the organisers'
lexical baseline. Token-overlap similarity reaches \JaccardTestAcc\% on test
against the \PublishedJaccardTest\% reported in the overview --- a difference of
\JaccardGapPts\ points, or \JaccardGapItems\ items out of \TestN.

The cause is tokenization. The pattern used here, \texttt{[a-z0-9']+}, treats
the apostrophe as part of the word, so \emph{father's} and \emph{doesn't} stay
single tokens; a plain \texttt{\textbackslash w+} pattern breaks them in two.
Across the test split there are \JaccardDiffTypes\ token types the first
pattern produces and the second never does, \JaccardDiffApostPct\% of them
apostrophe-bearing and the rest mostly fragments of accented words. Because the
score compares token sets, the two patterns pick different candidates on
\JaccardFlipN\ of the \TestN\ triples: on \JaccardFlipFor\ the pattern used
here picks the candidate the annotators chose, and on \JaccardFlipAgainst\ the
plain pattern does. That leaves the pattern used here \JaccardGapItems\ items
ahead --- \JaccardGapPts\ points on \TestN\ triples, which is the whole of the
gap. Four ordinary alternatives reproduce the published
\PublishedJaccardTest\% exactly. I retain the original tokenizer for
consistency with the rest of the pipeline, noting that it flatters this
baseline. No comparison here depends on it, and the development result of
Section~\ref{sec:analysis} is unaffected: all \JaccardVariantsN\ tokenizations
fall below chance there, between \JaccardMinVariant\% and
\JaccardMaxVariant\%.

% ===========================================================================
\section{Approach}\label{sec:approach}
% RUBRIC 3 (22 pts). What was built, why it was expected to work, and enough
% detail to reimplement. NO ACCURACIES IN THIS SECTION.
% ===========================================================================
The overview's catalogue of \OfficialRankedA\ ranked Track~A systems shows that
nobody used semantic role labelling and nobody used coreference resolution
\citep{hatzel2026} --- the two standard tools for recovering who acts and who is
affected --- while the dataset's own annotators record course of action as
driving a large share of decisions. The hypothesis here is therefore a gap
rather than a hunch: the field reached for large language models first, and the
obvious structural approach was never tested. What follows is an incremental
ablation in which each layer adds one representation, so that each
contribution is isolated rather than inferred.

\textbf{L0} is the floor: seeded random choice, and token-set Jaccard over
\texttt{[a-z0-9']+} on lowercased text. \textbf{L1} is TF--IDF cosine
(scikit-learn, bigrams, \texttt{min\_df}~$=2$, sublinear term frequency,
English stopwords), which asks how much of the task is simply word overlap; if
the answer were most of it, nothing structural would be needed. \textbf{L2} is
sentence embeddings, \texttt{intfloat/e5-base-v2} \citep{e5} with the
\texttt{"query: "} prefix, encoding whole summaries and ranking by cosine.
Sentence embeddings should help if the decision turns on paraphrase rather than
on structure. The organisers publish no Track~A figure for this encoder, so L2
also supplies a reference point the rest of the stack is measured against.

\textbf{L3} adds characters. Coreference chains from
\texttt{maverick-mes-litbank} \citep{maverick,jm3} give each story its cast, and
story similarity is the score of the best one-to-one matching between the two
casts, computed with the Hungarian algorithm
(\texttt{scipy.optimize.linear\_sum\_assignment}) over mention-string and
context similarity. For this to help, similar stories must be populated by
correspondable casts --- precisely what is being tested.

\textbf{L4} adds roles. For each character the verb lemmas where it is the
syntactic subject act as an agent proxy and those where it is the object as a
patient proxy \citep[Ch.~22]{jm3}, and the matching scores shared role--lemma
pairs instead of mentions. Where L2 scores what the story says and L3 who is in
it, L4 scores what is done: stories holding the same activities can land
together no matter who performed them. The proxy is deliberate rather than convenient, because three routes to
predicate--argument structure were available and none is clean. A trained
PropBank system would be the standard choice, but no maintained one exists;
AllenNLP's was archived in 2022. Prompting a language model is the current
substitute, and it is weak, reaching \SrlArgMatchLLM\% argument match against
\SrlArgMatchHuman\% for humans \citep{bonn2024}; it would also place inside the
representation the very class of model the structural hypothesis is meant to be
tested against. The third route, taken here, reads roles off the dependency
parse, as \citet{chaturvedi2018} did on this genre. Its cost is that
dependency relations diverge from semantic roles systematically: \emph{John broke the window} and \emph{the window broke} assign
the same participant different relations.

Because coreference is expensive and itself error-prone, L3 and L4 each have a
\textbf{coreference-free variant} in which the cast is built from spaCy
(\texttt{en\_core\_web\_sm}) noun-chunk heads. Two switches then vary
independently: an \texttt{actors\_only} filter, which keeps proper nouns and
nouns appearing as the subject of some verb; and span embedding with
\texttt{all-MiniLM-L6-v2}, which makes the alignment compare meanings rather
than strings. All four combinations were run on development data and, with the
coreference path, isolate the contributions of coreference and of cast purity.
The configuration carried to the held-out test is the one that scored best
there: noun-chunk cast, filter on, spans embedded, denoted L3x below.

\textbf{L5} combines the L2, L3 and L4 scores in a weighted sum with weights
fitted by cross-validation: if the layers represent different things, at least
one should correct the others' mistakes.

\textbf{L6} is sequential structure. Each summary is cut into five ordered
segments, each is embedded with \texttt{all-MiniLM-L6-v2}, and the stories are
compared slot by slot, because a whole-summary vector cannot distinguish a
story that resolves early from one that resolves late. A second variant relaxes
the correspondence by dynamic time warping, so the difference between the two
measures pacing rather than order. Five segments matches the story grammar that
received weight \GrammarWeight\ in a competing system's learned combination
here, against \TfidfWeight\ for TF--IDF and \EmbedWeight\ for embeddings.
\textbf{L6 was implemented but never evaluated}, and
Section~\ref{sec:ethics} returns to it.

\textbf{L7} decomposes each summary along the three constructs annotators
recorded about their own decisions --- course of action, outcomes and abstract
theme --- and scores the candidates construct by construct. If the annotation
reflects how the judgement is actually made, scoring them separately ought to
recover it.

\subsection{Harness}
Every layer implements one interface, \texttt{score(anchor, candidate)},
selected by Hydra configuration, so comparing two layers changes exactly one
thing. A single evaluation path serves all of them, giving bootstrap and Wilson
intervals for one system, McNemar's paired test for two \citep{dietterich},
cross-validation wherever parameters are fitted, and a guard refusing the test
split unless unlocked. Every figure here is a macro generated from the stored
results, with a provenance file recording each source and a lint that fails the
build on any hand-typed number. This caught three errors that would otherwise have reached the page. Two
configurations sharing a filename were silently overwriting each other's
results, so every result file now carries a run-identity hash; a
TF--IDF vectoriser was being fitted on the split it was scored on, and cost a
point once corrected; and the baseline discrepancy of Section~\ref{sec:data}
surfaced only because the audit refused an unreproduced published figure.


% ===========================================================================
\section{Results}\label{sec:results}
% RUBRIC 4 (17 pts).
% TODO CUT: this section plus Section 5 must fit roughly 2,000 words together.
% ===========================================================================

% CLAIM    One jump, then a plateau.
% BUDGET   ~600 words plus the table. Do not narrate the table row by row;
%          state the jump, state the plateau, then the paired tests.
% ---------------------------------------------------------------------------

\begin{table}[htb]
\centering\scriptsize
\setlength{\tabcolsep}{2.5pt}
\begin{tabular}{lrrr}
\toprule
\textbf{Layer} & \textbf{Acc.} & \textbf{SD} & \textbf{95\% CI} \\  % ok
\midrule
L0 \ random                    & \RandomAcc        & \RandomSD        & \RandomCI \\
L0 \ Jaccard                   & \JaccardAcc       & \JaccardSD       & \JaccardCI \\
L1 \ TF--IDF cosine            & \TfidfAcc         & \TfidfSD         & \TfidfCI \\
L2 \ \texttt{e5} embeddings    & \textbf{\LTwoAcc} & \LTwoSD          & \LTwoCI \\
\midrule
L3n cast, all nouns            & \LThreeNAcc       & \LThreeNSD       & \LThreeNCI \\
L3n cast, actors               & \LThreeNActorsAcc & \LThreeNActorsSD & \LThreeNActorsCI \\
L3x cast, emb., all nouns & \LThreeXAcc       & \LThreeXSD       & \LThreeXCI \\
L3x cast, emb., actors         & \LThreeXActorsAcc & \LThreeXActorsSD & \LThreeXActorsCI \\
L3x cast, emb., \texttt{e5} & \LThreeXEFiveAcc  & \LThreeXEFiveSD  & \LThreeXEFiveCI \\
\midrule
L4n roles, all nouns           & \LFourNAcc        & \LFourNSD        & \LFourNCI \\
L4n roles, actors              & \LFourNActorsAcc  & \LFourNActorsSD  & \LFourNActorsCI \\
L4x roles, emb., all nouns & \LFourXAcc        & \LFourXSD        & \LFourXCI \\
L4x roles, emb., actors        & \LFourXActorsAcc  & \LFourXActorsSD  & \LFourXActorsCI \\
\midrule
L5 \ combination, CV           & \LFiveXAcc        & \LFiveXSD        & \LFiveXCI \\
\bottomrule
\end{tabular}
\caption{Development results, $n=\DevN$. The \textbf{n} variants build the cast
without coreference; \textbf{x} variants embed the spans. SD is the spread of
individual items, which for a binary outcome follows from the accuracy, so the
interval is what describes uncertainty.}
\label{tab:dev}
\end{table}
\subsection{The ablation separates nothing}
The largest gap anywhere in the stack runs from lexical overlap to sentence
embeddings, \McnTfidfVsLTwoDiff\ points. Even that one is only suggestive:
McNemar places it at $p=\McnTfidfVsLTwoP$ over \McnTfidfVsLTwoDisc\ discordant
pairs. The backbone and the best structural layer are closer still. They
disagree on \McnLTwoVsLThreeXDisc\ of the \DevN\ items, and those disagreements
split \McnLTwoVsLThreeXB\ to \McnLTwoVsLThreeXC, so neither layer is ahead of
the other ($p=\McnLTwoVsLThreeXP$). The same holds for every other pair in the
table: no comparison on development data reaches significance. This is what the
power analysis of Section~\ref{sec:task} predicted in advance, and it is a
property of the sample size rather than a disappointment. At \DevN\ items
accuracy simply cannot adjudicate between these representations, which is why
Section~\ref{sec:analysis} does not ask it to.

\subsection{Held-out test result}

% CLAIM    The confirmatory claim, made once, on the frozen configuration.
% EVIDENCE Fill in after the single test run. Nothing in this section may be
%          written before that run, and the configuration may not change
%          after it.
% BUDGET   ~250 words.
% ---------------------------------------------------------------------------
The configuration was frozen at commit \texttt{\FreezeSha} in the public
repository,\footnote{\href{https://github.com/kclaudeeager/semeval2026-task4}{github.com/kclaudeeager/semeval2026-task4}}
the layers and the single confirmatory comparison fixed beforehand, and the
split evaluated once. Later commits change analysis and reporting only, never a
scorer or an evaluation path.

\begin{table}[htb]
\centering\footnotesize
\setlength{\tabcolsep}{3pt}
\begin{tabular}{lrr}
\toprule
\textbf{Layer} & \textbf{Test} & \textbf{Dev} \\
\midrule
L2 \ \texttt{e5} embeddings        & \textbf{\LTwoTestAcc} & \LTwoAcc \\
L5 \ combination, coref. comp.     & \LFiveCorefTestAcc    & --- \\
L0 \ Jaccard                       & \JaccardTestAcc       & \JaccardAcc \\
L5 \ combination, pre-reg. comp.   & \LFiveTestAcc         & \LFiveXAcc \\
L1 \ TF--IDF                       & \TfidfTestAcc         & \TfidfAcc \\
L3x character alignment            & \LThreeXTestAcc       & \LThreeXActorsAcc \\
L0 \ random                        & \RandomTestAcc        & \RandomAcc \\
\bottomrule
\end{tabular}
\caption{The held-out run, $n=\TestN$, beside the development figures. Every
layer improves on held-out data except character alignment, which is the only
one that falls.}
\label{tab:test}
\end{table}

The confirmatory claim is supported, more strongly than development data
indicated. The backbone reaches \LTwoTestAcc\% against \LThreeXTestAcc\% for
the best structural layer, a difference McNemar places at
$p = \McnLTwoTestVsLThreeXTestP$ over \McnLTwoTestVsLThreeXTestDisc\ discordant
pairs; the pre-registered combination reaches \LFiveTestAcc\%, also below the
backbone ($p = \McnLTwoTestVsLFiveTestP$). On development data no layer
separated from any other, whereas on held-out data the backbone separates from
both.

The only combination matching the backbone was the one whose fitted weights
discarded the structural layers entirely (\LFiveCorefWeights).

The pre-specified exploratory measure did not replicate, which is the outcome
the pre-registration named in advance as falsifying it. Cross-entropy under the
anchor's unigram model selected the distractor on development data at
$p=\LexContentCrossEntP$, but favours the correct candidate on
\LexTestCrossEntPct\% of test triples ($p=\LexTestCrossEntP$), an interval
containing chance. The lexical reversal therefore belongs to the development
split alone. One further pattern is visible in Table~\ref{tab:test}: every
layer scores higher on test than on development except character alignment,
which falls. Having tied the backbone on development, it is six points behind
on held-out data, so that parity was itself an artifact.

A deviation from the pre-registration is recorded there and repeated here: it
named the combination over the embedded-span variants, but the first run used
the configuration default, which is the coreference set. The pre-registered
variant was run afterwards and both appear above; neither is withheld.

% ---------------------------------------------------------------------------

% ===========================================================================
\section{Analysis}\label{sec:analysis}
% RUBRIC 5 (27 pts). Concrete examples are worth 16 of the 27.
% ===========================================================================
Accuracy cannot adjudicate at \DevN\ items, so the question is attacked by
other means. Two checks bound what any alignment layer could have achieved. The
first is geometric: placing the anchor at the origin, the chosen candidate is
the nearer of the two in \GeomTfidfCloserPct\% of triples in TF--IDF space and
\GeomMiniLMCloserPct\% in MiniLM space, neither separable from a coin flip, and
reliably nearer only in \texttt{e5} space ($p=\GeomEFiveCloserP$). Since the
cost matrices of Section~\ref{sec:approach} are built in MiniLM space, the one
space that carries signal is the one the alignment never enters; and because
that signal is whole-story, decomposing a summary into spans discards precisely
what \texttt{e5} encodes. The second check concerns the annotated dimensions:
scored alone over \CompN\ triples, theme reaches \CompThemeAcc\%, course of
action \CompActionAcc\% and outcomes \CompOutcomeAcc\%, every interval
containing chance, while weighting all three reaches only \LSevenAcc\%, below
the undecomposed backbone at \LTwoAcc\%.

\subsection{Surface overlap is reversed on development data}
Before any comparative measurement of meaning, one would try the simplest text
comparison available: the ratio of words the anchor and a candidate share. It
is not a smart instrument, since synonyms and argument structure are invisible
to it, but it is worth checking before reaching for anything more elaborate.

On development data the distractor's overlap with the anchor is \emph{higher}
than the right answer's --- accuracy \JaccardAcc\% with a mean signed margin of
\JaccardMargin. A distributional measure confirms it: cross-entropy of the
candidate under the anchor's unigram model favours the correct candidate on only
\LexContentCrossEntPct\% of triples ($p = \LexContentCrossEntP$).

One triple shows the mechanism. The anchor concerns a biographer investigating
a fatal test flight, who falls in love with the dead pilot's sister as people
connected to the case begin to die. The distractor describes a chemist
betrothed to one sister and in love with another, extorted by agents over his
secret work; because it shares the sibling, the romance, the state secrets and
the register, it reaches a Jaccard of \ExampleJaccardWrong. The candidate the
annotators chose shares far less vocabulary, \ExampleJaccardRight, describing a
teenager whose family is murdered by foreign agents and who discovers her life
was a fabrication. What it does share is the \emph{shape} of the anchor, in
which a protagonist uncovers a concealed truth and is endangered by it. Overlap
therefore prefers the distractor, by \ExampleJaccardGap.

This is a property of the construction rather than a usable signal. The
distractors were drawn at the \DistractorPercentile\ percentile of an encoder
that rewards shared surface, so the effect is built in; but it cannot be
exploited, because the rule does not hold on unseen data, where the same
baseline reaches \JaccardTestAcc\%. What overlap cannot see is argument
structure, and that is what motivates the layers below.

\subsection{Identifying characters more carefully does not help}
The next intuition is to ask who the story is about, since the matching story
should share more characters with the anchor than the distractor does. The cast
extraction tests that premise directly. A noun-chunk head is kept when it is a
proper noun, or when that noun appears as the subject of a verb somewhere in
the text, which makes the filter behavioural rather than categorial: it asks
whether a thing ever acts, not whether it is a person.

On one development anchor, about an organisation advocating for future
generations, the filter reduces a cast of \ExampleCastFull\ noun heads to
\ExampleCastActors. The \ExampleCastRemoved\ discarded include
\texttt{citizen}, \texttt{minister}, \texttt{advocate}, \texttt{glacier} and
\texttt{airship} --- settings, objects and abstractions. If the alignment were
matching characters, discarding these should help.

It does not. Restricting the cast costs \McnLThreeNVsLThreeNActorsDiff\ points
for the character layer, from \LThreeNAcc\% to \LThreeNActorsAcc\%, and
\McnLFourNVsLFourNActorsDiff\ for the role layer. Proper coreference costs more still: those
variants reach \LThreeAcc\% and \LFourAcc\%, both intervals containing
chance. None of the differences is significant
($p = \McnLThreeNVsLThreeNActorsP$, $p = \McnLFourNVsLFourNActorsP$), so the
argument rests on direction rather than margin. What makes that direction hard to dismiss is its
consistency: four comparisons, across two independently built layers and two
mechanisms for identifying characters, all move the same way.

The direction inverts the expected reading. Discarding non-character nouns made
the system worse, so those nouns were contributing to correct decisions:
whatever the alignment matched on, it was not characters specifically, and much
of the cost matrix was scoring settings, objects and topic words --- the same
vocabulary overlap already shown to be reversed on this split. One caveat: the
coreference path builds its cast from resolved chains and the coreference-free
path from noun chunks, so that comparison changes the mechanism of cast
construction rather than isolating coreference.

\subsection{The structural layers err like the lexical baseline}
If a structural representation sees something word overlap cannot, it should
fail on \emph{different} items. Correlating each layer's per-item errors with
the TF--IDF baseline's over the same \DevN\ triples tests that without
reference to accuracy at all.

\begin{table}[h]
\centering\small
\begin{tabular}{lr}
\toprule
\textbf{Layer} & \textbf{$\phi$ with TF--IDF} \\
\midrule
L2 \quad sentence embeddings alone    & \PhiLTwo \\
L5 \quad combination, cross-validated & \PhiLFiveX \\
L3x \quad characters, embedded spans  & \PhiLThreeX \\
L4x \quad roles, embedded spans       & \PhiLFourX \\
L3n \quad characters, overlap inside  & \PhiLThreeN \\
L4n \quad roles, overlap inside       & \PhiLFourN \\
\bottomrule
\end{tabular}
\caption{Correlation between each layer's errors and the lexical baseline's, on
\DevN\ development triples. Lower means the layer fails on different items,
which is what a different representation should do.}
\label{tab:phi}
\end{table}

Every value is positive and none attributable to chance: shuffling a layer's
error vector while holding accuracy fixed places $\phi$ between \PhiNullLo\ and
\PhiNullHi\ ($p=\PhiNullPLTwo$ for the embedding layer, $p=\PhiNullPLThreeX$
for character alignment). The prediction therefore fails, and it fails in a
particular direction, since sentence embeddings alone are the \emph{least}
lexically correlated layer at \PhiLTwo\ and every structural layer sits above
them. The ordering tracks how much surface matching remains inside each one.
The combination is the clearest case: at \PhiLFiveX\ it sits above the backbone
it was built on, so adding structure moved its errors \emph{closer} to the
lexical baseline. Shared item difficulty would raise every coefficient equally
and so cannot explain why they differ, which is why the ordering is the finding
rather than the magnitudes.

\subsection{The combination ceiling is far below independence}
Layers that fail together cannot be combined into something better than
themselves. The union, the proportion of items on which at least one layer is
correct, is not itself a system, since selecting the right layer per item would
presume knowing the answer; but it is the ceiling for every combination rule. Over the \PoolSize\ layers pooled here at least one is
correct on \OracleUnion\% of development items, where independence would
predict \OracleIndep\%. That shortfall of \OracleShortfall\ points is the
error correlation above, expressed in accuracy. Both realised combinations land
\emph{below} the backbone, so combining did not merely fail to help but cost.

% ===========================================================================
\section{Ethics and Limitations}\label{sec:ethics}
% RUBRIC 6 (10 pts).
% ===========================================================================
\paragraph{Licensing.} The summaries are licensed CC BY-SA 4.0, the labels  % ok
CC0. Summaries may therefore be redistributed with attribution and under the
same licence, which this report observes in quoting them, while the labels
carry no restriction. The underlying texts are encyclopaedic plot summaries of
published films and novels, so the corpus remains derivative of copyrighted
works even though the summaries themselves are openly licensed; the examples in
Section~\ref{sec:analysis} are accordingly excerpted for analysis rather than
reproduced. Finally, the dataset README asks that a canary string accompany any
reproduction of samples with their labels, and the script that generates the
examples below emits it alongside them.

\paragraph{Test-label discipline.} The gold test labels are public, so using
them is a matter of restraint rather than of access, and one team in this
shared task was disqualified for failing that restraint \citep{hatzel2026}.
Three mechanisms were used in place of good intentions. The evaluation code
refuses the test split unless an environment flag is set. The confirmatory
comparison, the layers to run and the exploratory measures were all written
down before the split was touched. And deviations found afterwards were
recorded in a dated section of that document rather than by editing it. The
reconciliation in Section~\ref{sec:data} was the one additional computation
performed on test, and it fitted and selected nothing.

\paragraph{Scope.} The low per-dimension agreement in Section~\ref{sec:data} is
a property of the task rather than evidence of careless annotation, since
people who agree on which story is closer can reasonably disagree about why.
The data is English only and drawn from a largely Western film and novel canon,
so a claim here about ``narrative similarity'' is a claim about that canon.
That matters for the interest motivating this project, structural
representations for low-resource storytelling: a negative result on compressed
English summaries is not decisive for genres where structure is explicit.

\paragraph{Use of generative AI.} I used Claude throughout this project. It wrote most of the supporting code ---
the evaluation harness additions, the audit script, and the figure and
provenance tooling --- proposed several of the diagnostic measurements, and
drafted report prose from my explanations, which I reviewed and revised. The
boundaries were set deliberately: the research question, the layer design,
every experiment run and every decision about what could honestly be claimed
were mine; no figure here comes from a model rather than from a stored result;
and the errors it introduced, which I caught, are part of why the verification
machinery of Section~\ref{sec:approach} exists.

\paragraph{Limitations.} The paired test on \TestN\ items detects
\PowerPairedLo\ to \PowerPairedHi\ points, so a small true effect would be
invisible, and stratified contrasts need \PowerStrataMDE\ points and are
therefore description only. Dependency relations proxy semantic roles and
diverge from them systematically, and the coreference comparison is confounded
with cast construction, since one path uses chains and the other noun chunks.
L0 and L1 measure surface-form overlap rather than lexical semantics, and the
stopword list in the distributional analysis is ad hoc. Stories recur across
splits, and only one encoder family and one language were used.

Two gaps deserve naming rather than burying. The first is that \textbf{L6 was
never evaluated}, even though it targets the one structural signal with a
positive published result on this task; it also inherits a mismatch, since its
segment count came from a \emph{functional} story grammar but is implemented as
\emph{positional} slicing, which on summaries this short means a segment of
about one sentence and \SegPadPct\% of summaries padded by a repetition. The
second is that, measured against the proposal, L7 and the geometric analysis
were unplanned additions and a promised hand-audit of coreference outputs did
not happen.

\paragraph{Conclusion.} Structure did not help, and the reason is measurable:
the alignment compares spans inside a space where even whole summaries fail to
separate the candidates, so decomposing a summary discards the signal the
encoder carries. The next step is to score that decomposition, not match it.

\clearpage
\bibliographystyle{acl_natbib}
\bibliography{custom}

\end{document}
